\section{Gauge transport and its continuum comparison}
Endpoint cancellation gives covariance of composed transport and conjugation of closed holonomy. In the specified finite Maxwell--Dirac system, matter continuity, Gauss propagation, and electric work fit together exactly. Ordered non-Abelian products retain the path information required by refinement.

Geometric interpolation supplies actual midpoints and weighted isometries. Straight-edge reconstruction and path curvature enter distinct error budgets for explicitly constructed finite operators. Regular line-integral links reconstruct straight edges exactly, while derivative envelopes give summable dyadic bounds for the chosen curvature defect.

Correctly normalized finite fields become $L^2$ step functions. Concrete dyadic boxes identify the approximating subspaces modulo null boundaries, and their projections converge strongly on each fixed field. Comparison of physical generators and evolution remains a separate task.

\subsection{The remaining analytic bridge}
The next background-coupled result needs a common operator setting. One must specify the physical connection, its regularity, the continuum spinor domain, the discrete operators, and the embeddings that compare them. Consistency on a controlled core, stability, and a suitable graph or resolvent argument then have to connect the finite refinement estimates to the continuum operator. Each limit must identify its topology and its dependence on the state or its regularity.

The free Dirac route already provides a useful endpoint for this comparison. The free construction gives a self-adjoint physical torus Hamiltonian on its weighted Fourier domain. The generator chapter supplies its spectral-evolution identification as a written proof. Its formal certification remains distinct from the recorded operator certificate. A background gauge extension must account for those operator-domain and evolution questions as well as the link estimates developed here.

\subsection{From a prescribed connection to coupled dynamics}
A background estimate controls a supplied field. A coupled gauge theory makes that field part of the unknown. The finite semidiscrete example illustrates the desired compatibility among matter transport, electric evolution, and a Gauss constraint. Extending this compatibility to a specified continuum setting requires an existence theory, constraint propagation, and a convergence statement for the coupled system. Local covariance alone supplies none of these additional conclusions.

The constrained-gravity chapter examines the additional obligations imposed by a dynamical geometry: regular-chart reduction, constraint closure, boundary terms, and strong-field evolution. Arithmetic refinement and cosmological transport retain the same requirement to identify their state spaces and comparison maps, without sharing all of the gauge model's assumptions.

\begin{readerroadmap}[title={The next theorem should name its comparison}]
Specify the two state spaces, the operators or evolutions to be compared, the map between resolutions, and the norm or topology in which the defect tends to zero. Then state which connection and domain hypotheses make that comparison valid.
\end{readerroadmap}

